\subsection{Outlook: the aggregated setting}
\label{sec:aggregated}

The algebraic setting cannot see analytic phenomena such as the Stone--von Neumann theorem, so it
is natural to ask which algebraic Witt classes become trivial once analytic objects are allowed. We
briefly describe one setting in which to ask this, and what can be said formally. In this section
$k=\CC$.

Aggregated vector spaces form a Grothendieck abelian category containing all complete bornological
spaces, e.g.\ Fr\'echet spaces, with a symmetric monoidal structure extending the completed
projective tensor product \cite{Steff:six, Schneiders}. Let $\Pre_{\Agg}$ be the $\infty$-category
of presentable $\infty$-categories enriched in its derived $\infty$-category, base changed to
$\CC$. A version $\QC^{\Agg}$ of $\QC$ on finitely generated affine derived $C^\infty$-schemes,
which include all compact manifolds and $\RR^n$, with a six-functor formalism satisfying K\"unneth
and base change, is constructed in \cite{Steff:six}; for compact $X$ it is aggregated modules over
$C^\infty\left(X\right)$. Base change $\Pre\to\Pre_{\Agg}$ is symmetric monoidal and preserves
colimits, so it induces a homomorphism of Witt groups (Definition \ref{defn:Witt})
\begin{equation}
\label{eqn:Witt_Agg}
\varphi\colon\mathrm{W}\left(\Pre\right)\longrightarrow\mathrm{W}\left(\Pre_{\Agg}\right) \ .
\end{equation}

For a compact torus $T$ (not the algebraic torus of Section \ref{sec:tori}), write
$\mathrm{Rep}\left(T\right)$ for the comodules over
$C^\infty\left(T\right)$ in aggregated modules, using
$C^\infty\left(T\right)\hat\otimes C^\infty\left(T\right)\simeq C^\infty\left(T\times T\right)$.

\begin{lem}
\label{lem:averaging}
$C^\infty\left(T\right)$ is descendable in $\mathrm{Rep}\left(T\right)$.
\end{lem}

\begin{proof}
Integrating against the normalized Haar measure defines a bounded map $C^\infty\left(T\right)\to\CC$,
and by translation invariance it is a map of comodules retracting the unit. So the unit is a retract
of $C^\infty\left(T\right)$, and in particular lies in the thick tensor ideal it generates
\cite{Mathew}.
\end{proof}

For a finite abelian group $\Gamma$ with a non-degenerate quadratic form $q$, we regard
$\Vect_\Gamma$ with the braiding determined by $q$ \cite{JS:braided}; it is invertible \cite{BJSS}
(Remark \ref{rmk:abelian}). Its central charge $c\in\ZZ/8$ is defined by
$\left|\Gamma\right|^{-1/2}\sum_{x\in\Gamma}q\left(x\right)=\exp\left(2\pi ic/8\right)$
\cite[Appendix 4]{MH}.

\begin{prop}
\label{prop:Witt_reduction}
Suppose given classes $a_V\in\mathrm{W}\left(\Pre_{\Agg}\right)$, one for each non-degenerate real
quadratic space $\left(V,b\right)$, such that
\begin{enumerate}
\item $a_V$ only depends on $\left(V,b\right)$ up to isometry;
\item $a_{V\dsum V'}=a_V+a_{V'}$;
\item for every even lattice $\left(N,b\right)$ with discriminant form $\left(\Gamma,q\right)$,
$\varphi\left[\Vect_\Gamma\right]+a_{N\otimes\RR}=0$.
\end{enumerate}
Then there is a class $g\in\mathrm{W}\left(\Pre_{\Agg}\right)$ with $8g=0$ such that
$\varphi\left[\Vect_\Gamma\right]=c\,g$ for every finite non-degenerate quadratic form
$\left(\Gamma,q\right)$ of central charge $c$. In particular, every such $\Vect_\Gamma$ with $c=0$
lies in the kernel of $\varphi$.
\end{prop}

\begin{proof}
Every $\left(\Gamma,q\right)$ is the discriminant form of an even lattice $\left(N,b\right)$
\cite[Theorem 1.10.1]{Nik}, so $\varphi\left[\Vect_\Gamma\right]=-a_{N\otimes\RR}$ by (3). Let
$h_\pm\ceqq a_V$ for $V=\RR$ with the form $\pm xy$. By Sylvester's law of inertia and (1), (2),
$a_V=r\,h_++s\,h_-$ for $V$ of signature $\left(r,s\right)$. The hyperbolic lattice is even with
trivial discriminant form, so (3) gives $h_++h_-=0$, and hence $a_V=\sigma\left(V\right)h_+$. The
$E_8$ lattice is even with trivial discriminant form and signature $8$, so $8h_+=0$. Finally
$\sigma\left(N\right)\equiv c$ modulo $8$ by Milgram's formula \cite[Appendix 4]{MH}, so
$g\ceqq-h_+$ works.
\end{proof}

So the relation $8g=0$ is forced by lattices, not by analysis. We expect such classes to arise as
$a_V=\left[\cA_V\right]$ for a braided category $\cA_V$ of sheaves on $V$ with braiding
$\exp\left(-2\pi i\,b\left(x,y\right)\right)$. Namely, for an even lattice $\left(N,b\right)$ and
$T=\left(N\otimes\RR\right)/N$, the form $b$ twists convolution on $\QC^{\Agg}\left(T\right)$ via a
biextension of $T\times T$, which has no algebraic counterpart since line bundles on algebraic tori
are trivial. Using Lemma \ref{lem:averaging}, the resulting $\EE_1$-algebra should Morita trivialize
its Drinfeld center $\cZ_N$, and a Peter--Weyl theorem for $T$, via Fourier transform, should
identify $\cZ_N$ with $\Vect_\Gamma\otimes\cA_{N\otimes\RR}$. This would give (3). Making this
precise requires defining $\cA_V$, including its decay conditions, and descent for module
categories along the descendable algebra of Lemma \ref{lem:averaging}.

\begin{question}
\label{qn:order_g}
Do such classes exist, and if so, what is the order of $g$? If $g=0$, which would be a $2$-shifted
analogue of the Stone--von Neumann theorem, every non-degenerate pointed braided fusion category lies in the kernel
of $\varphi$. Otherwise $\varphi$ detects part of the central charge modulo $8$.
\end{question}
